\documentclass[12pt,a4paper]{article}
\usepackage[margin=1in]{geometry}
\usepackage[utf8]{inputenc}
\usepackage[T1]{fontenc}
\usepackage{lmodern}
\usepackage{microtype}
\usepackage{parskip}
\usepackage[colorlinks=true,urlcolor=blue]{hyperref}
\setlength{\parindent}{0pt}
\setlength{\parskip}{0.9em}

\begin{document}

\begin{flushright}
Wang Weijia\\
School of Economics and Finance\\
Xi'an Jiaotong University\\
Xi'an, Shaanxi, China\\[0.3em]
\texttt{18522704823@163.com}\\[0.5em]
\today
\end{flushright}

\bigskip

Dear Guest Editors,

I am pleased to submit the manuscript ``\textbf{Artificial Intelligence, Skill Mismatch, and Employment Reallocation in China}'' for consideration in \textit{China Economic Review}'s special issue \textit{Income and Wealth Distribution, Labor Market, and Subjective Well-being in China---In Honor of Professor John Knight}.

\textbf{Summary of contribution.} This paper estimates how firm-level AI adoption affects employment quantity and skill composition in Chinese A-share listed firms from 2010 to 2022, and how pre-existing overeducation and skill mismatch condition these effects. Using an education-expansion shift-share instrument (Edushock) and a leave-one-out peer AI diffusion instrument, I find that AI adoption raises total employment (IV coefficient: 0.574) and increases the high-skilled labor share by 0.109 while reducing the low-skilled share by 0.150. Overeducation weakens AI's complementarity with high-skilled workers but cushions low-skilled displacement. Within-firm pay dispersion data provide direct evidence that AI adoption is associated with a wider executive-to-worker compensation gap (IV coefficient: 0.031), an effect that overeducation statistically compresses. These employment and wage findings are derived from and interpreted within the Acemoglu--Restrepo task model, which provides precise predictions linking AI-driven task reallocation to the skill wage premium.

\textbf{Connection to the special issue.} The paper speaks directly to the special issue's central themes. Professor Knight's foundational work on Chinese wage inequality and labor-market adjustment provides the empirical and conceptual anchor for this paper's income distribution analysis. The results show that the current wave of AI adoption continues and potentially amplifies the pattern of widening skill-wage differentials documented in Knight and Song (2003) and Knight (2014): AI expands employment but concentrates job creation and pay gains among high-skilled workers, while overeducation---the mismatch accumulated during China's rapid credential expansion---moderates the magnitude of this distributional shift. The findings therefore contribute to understanding how a new technology shock interacts with structural features of China's labor market to shape income distribution, a question at the core of Professor Knight's research agenda.

\textbf{Originality and related work.} This paper has not been published, accepted for publication, or submitted for review elsewhere. It is original work and is not under consideration at any other journal. All empirical results are based on the author's own data construction and estimation.

\textbf{Declarations.} The author declares no competing financial interests or personal relationships that could have influenced this work. No external funding was received for this research.

I would welcome the opportunity to have the paper reviewed by referees with expertise in Chinese labor economics, automation and the labor market, and IV identification strategies. I am available to provide any supplementary materials the editorial team requires.

Thank you for considering this submission.

\bigskip

Sincerely,

\bigskip\bigskip

Wang Weijia\\
School of Economics and Finance, Xi'an Jiaotong University\\
\texttt{18522704823@163.com}

\end{document}
